\documentclass[12pt,a4paper]{article}
\usepackage[utf8]{inputenc}
\usepackage[spanish]{babel}
\usepackage{amsmath,amssymb,amsfonts}
\usepackage{geometry}
\usepackage{hyperref}
\geometry{margin=2.5cm}

\title{\textbf{Constantes Matemáticas en Cálculo Vectorial}}
\author{Compendio de Constantes Fundamentales}
\date{\today}

\begin{document}

\maketitle

\section*{Introducción}
El cálculo vectorial extiende el cálculo diferencial e integral a campos escalares y vectoriales en espacios euclídeos multidimensionales ($\mathbb{R}^2, \mathbb{R}^3, \dots$). Las constantes asociadas al cálculo vectorial determinan medidas topológicas fundamentales como el área de superficies cerradas, volúmenes integrales y constantes de integración geométrica espacial.

\section{Área de la Esfera Unitaria ($A_{S^2} = 4\pi$)}

\subsection*{1. Significado, Símbolo y Explicación}
El área de la superficie de la esfera de radio unidad ($r=1$) en $\mathbb{R}^3$, denotada por $A_{S^2}$, se calcula mediante la integral de superficie cerrada:
\begin{equation*}
A_{S^2} = \iint_{S^2} dS = \int_{0}^{2\pi} \int_{0}^{\pi} \sin\theta \, d\theta \, d\phi = 4\pi
\end{equation*}
Representa la medida total del ángulo sólido que rodea completamente un punto en el espacio tridimensional ($4\pi$ estereorradianes) y aparece como factor de normalización fundamental en el Teorema de la Divergencia de Gauss y la ley de Coulomb.

\subsection*{2. Valor Típico en Ingeniería y Ciencia}
\begin{equation*}
4\pi \approx 12.5664 \quad \text{o} \quad 12.56637061
\end{equation*}

\subsection*{3. Valor Exacto a 15 Decimales}
\begin{equation*}
4\pi \approx 12.566370614359173
\end{equation*}

\subsection*{4. Valor Exacto a 100 Decimales}
\begin{align*}
4\pi \approx \; & 12.5663706143591729538505735331180115367886775975004232838997783692312... \\
& ...656251448359445130136848847231920860
\end{align*}
\textbf{Margen de error:} Constante matemática pura ($4 \times \pi$).

\vspace{0.8cm}
\hrule
\vspace{0.8cm}

\section{Volumen de la Bola Unitaria ($V_{B^3} = \frac{4}{3}\pi$)}

\subsection*{1. Significado, Símbolo y Explicación}
El volumen de la bola sólida de radio $r=1$ en $\mathbb{R}^3$, denotado por $V_{B^3}$, se obtiene integrando el elemento de volumen en coordenadas esféricas:
\begin{equation*}
V_{B^3} = \iiint_{B^3} dV = \int_{0}^{1} r^2 dr \int_{0}^{\pi} \sin\theta d\theta \int_{0}^{2\pi} d\phi = \frac{4}{3}\pi
\end{equation*}
Es la constante volumétrica fundamental en teoremas de integración vectorial, dinámica de fluidos y teoría del potencial.

\subsection*{2. Valor Típico en Ingeniería y Ciencia}
\begin{equation*}
\frac{4}{3}\pi \approx 4.1888 \quad \text{o} \quad 4.18879020
\end{equation*}

\subsection*{3. Valor Exacto a 15 Decimales}
\begin{equation*}
\frac{4}{3}\pi \approx 4.188790204786391
\end{equation*}

\subsection*{4. Valor Exacto a 100 Decimales}
\begin{align*}
\frac{4}{3}\pi \approx \; & 4.1887902047863909846168578443726705122628925325001410946332594564104... \\
& ...218750482786481710045616282410640287
\end{align*}
\textbf{Margen de error:} Constante matemática exacta (sin margen de error).

\end{document}
